% RESEARCH_DRAFT_SHA256: 6978402f264c8d660ccb1c69384a4c22fc9518ac405fc628295bc2ee07e8700c
% Layout and source-reference wrappers expanded for the web renderer.
\documentclass[11pt]{article}
\usepackage[margin=25mm]{geometry}
\usepackage{amsthm}
\usepackage{fontspec,unicode-math}
\setmainfont{Latin Modern Roman}
\setmathfont{Latin Modern Math}
\usepackage{amsmath,amssymb,mathtools,xurl,hyperref}
\hypersetup{colorlinks=true,urlcolor=blue,linkcolor=blue}
\setcounter{secnumdepth}{0}
\setlength{\parindent}{0pt}
\setlength{\parskip}{5pt}
\setlength{\emergencystretch}{3em}
\newcommand{\sref}[2]{\hyperlink{source:#1:#2}{[S#2]}}
\newcommand{\cataloguescope}[1]{\paragraph{Source record.}#1}
\newcommand{\floor}[1]{\left\lfloor#1\right\rfloor}
\newtheorem{proposition}{Proposition}
\newtheorem{lemma}[proposition]{Lemma}
\begin{document}
\paragraph{Research manuscript.} OpenAI \texttt{gpt-6-astra}, reasoning \texttt{ultra}; 11 October 2026.
\paragraph{Preserved source body.} The complete catalogue mathematical text and sources follow with layout and source links expanded, including source qualifications and its historical preparation notes. New research begins at the section headed Research attempt.

% UnsolvedMath ID 20002141; source catalogue code AIM-INFRASTRUCTURE-0041
\section{A certifiable registry profile for mathematical visualizations}\label{20002141}
{\small UnsolvedMath ID 20002141 · Computer Science}
\subsection{Definitions and mathematical statement}

\paragraph{Shared notation.}Standard number systems are $\mathbb N,\mathbb Z,\mathbb Q,\mathbb R,\mathbb C$. Specialized notation and index ranges follow the cited source and the definitions below.


A mathematical visualization represents mathematical objects, relations, or data visually. A catalogue can associate a visualization with its topic, explanatory text, source, license, and technical requirements. An interactive visualization changes in response to user input; a static visualization has a fixed presentation.
\paragraph{Record scope and qualifications.}
The catalogue title is retained as metadata. Named certificates, contracts, bounds or frameworks in that title are not asserted as results of the original source extract.
The extract permits outputs ranging from a plan to a prototype. It gives no fixed database schema, coverage requirement, or algorithmic target.

Catalogue screening limitation: Workshop planning or software infrastructure agenda; the supplied passage does not identify a determinate mathematical target.

The source snapshot labels this record partially solved; this is catalogue metadata rather than independent verification of a mathematical result or the current completion of a project.
\paragraph{Statement recorded in the supplied source.}
The goal of this group is to discuss ways to create an online database of mathematical visualizations. While creating such a database in a week is not feasible, the group aims to find tools that would make it possible and to show that such a project is possible in the first place. The output can be anything from a written plan of action, to an example of how it would work in the first place. \href{http://aimpl.org/cyberinfrastructure/4/}{[S2]}
\paragraph{Source status.}
The supplied catalogue labels this source record partially solved; that metadata was checked against the supplied snapshot on 2026-10-11. The label is preserved for provenance and is not an independent verification of present open status or mathematical correctness.
\subsection{Short English statement}
Plan and demonstrate infrastructure for an online database of mathematical visualizations.
\subsection{Sources}
\begin{itemize}
\item[S1] Ulam AI and the UnsolvedMath Contributors, UnsolvedMath: A Curated Collection of Open Mathematics Problems, record 20002141 (AIM-INFRASTRUCTURE-0041). Catalogue curation and metadata: CC BY 4.0. The source extract is retained; definitions, English summaries and qualifications are editorial additions.
\url{https://huggingface.co/datasets/ulamai/UnsolvedMath}.
\item[S2] AIM Workshop Problem Lists, source linked by catalogue record AIM-INFRASTRUCTURE-0041.
\url{http://aimpl.org/cyberinfrastructure/4/}.
\end{itemize}
\paragraph{Preparation and review.}
The original source and UnsolvedMath contributors retain attribution for the source extract. Codex (GPT-6; exact serving revision not exposed) assisted with definitions, English summaries, scope notes and formatting. No new human mathematical authorship or independently verified solution is asserted. Permanent source ID 20002141 and code AIM-INFRASTRUCTURE-0041 checked against the complete supplied registry and current repository. Distinct upstream records, including repeated questions and broad or damaged extracts, are retained for complete coverage of the source CS category. A repeated record is not a new independent conjecture. Complete mathematical-equivalence and current-literature review has not been performed. Preparation and final review follow contribution prompts C and D, with source gaps disclosed. Structural and rendering checks do not establish a new theorem or open status.
\label{end:20002141}

\section{Research attempt}
\subsection{Literature, source context and scope}
The complete catalogue body above retains the original agenda, source status, and historical preparation. The AIMPL problem page could not be retrieved. The supplied complete source transcription was read, and the accessible primary AIM workshop announcement was checked for context. The announcement is not a substitute for verification of the inaccessible problem page. The catalogue's earlier proposed frameworks and research summaries are context, not established solutions credited to this attempt. The following is an explicitly limited investigation of one route; it neither substitutes a formalized toy model for the entire source agenda nor claims empirical or deployed-software results. Elementary auxiliary facts are not new original-target progress.

\subsection{Testing a visualization certificate based on sampled agreement}
An online visualization collection needs to distinguish a displayed approximation from a certified mathematical object. Test the route of certifying a scalar graph by agreement on finitely many sample abscissas $x_1,\ldots,x_m$. Assume these are distinct real numbers.
\begin{proposition}
Sample agreement alone does not identify even a polynomial function. If instead a function is known to be $L$-Lipschitz on a sampled interval and adjacent sample gaps are at most $h$, its nearest-sample approximation has error at most $Lh/2$.
\end{proposition}
\begin{proof}
The polynomials $f(x)=0$ and $g(x)=\prod_{j=1}^m(x-x_j)^2$ agree at every sample, but $g(x)>0$ at every real $x$ outside the sample set. Multiplying $g$ by an arbitrary positive constant makes the discrepancy at any such point arbitrarily large. For the second claim assume the interval endpoints are included among the samples. Every point lies between two adjacent samples at distance at most $h$; its distance to the nearer endpoint is at most $h/2$. Lipschitz continuity gives $|f(x)-f(x_j)|\le L|x-x_j|\le Lh/2$ for that sample.
\end{proof}
Thus a finite screenshot or sample table by itself cannot certify all features between sampled points. A bound on regularity changes the available guarantee, but verifying the bound is a separate obligation. The result is restricted to one-dimensional scalar graphs; it does not automatically validate topological, geometric or interactive visualizations.
\subsection{Verification and remaining gaps}
For samples $0,1/4,1/2,3/4,1$, the polynomial $g$ was evaluated exactly at every sample and at $1/8$, where it is $11025/1073741824>0$. No visualization registry, search tool, metadata schema, or general certification method was built. The exact counterexample locates a failure of sample-only certification.

\section{References and source attribution}
Primary source: \url{http://aimpl.org/cyberinfrastructure/4/}. Attempted primary AIMPL page retrieval; failed with timeout or gateway error. The supplied full source transcription was read; the inaccessible live page was not independently inspected.

\url{https://aimath.org/pastworkshops/cyberinfrastructure.html}. Read the primary December 2023 workshop announcement; it describes the open-source mathematical infrastructure agenda.

All retrievals were checked on 11 October 2026. Original contributor and catalogue attribution remains in the preserved Sources section. The elementary arguments below are self-contained; their exposition is not a claim of first discovery.


\paragraph{Integration review.} Codex (GPT-6; exact serving revision not exposed) reviewed the written arguments and expanded layout and source-reference wrappers. Repository metadata and the production web renderer were checked using Python 3.13 and JavaScript with MathJax 3.2.2. This model-assisted review is not independent human proof verification; no PDF compilation or formal proof checking is asserted.

\section{Completion Estimate}
\noindent\textbf{COMPLETION: 0\%}
No substantiated new progress toward the entire original target is claimed. The worked results and failed route are delimited above; a known result, illustrative model, or auxiliary obstruction is not counted as original-target completion.

\section{Final status and attribution}
\textbf{Attempt status: unresolved.} This describes this research attempt, not a new ruling on the source catalogue's status or on external literature claims.
The mathematics and this self-audit were generated by OpenAI \texttt{gpt-6-astra}, reasoning \texttt{ultra}, on 11 October 2026. Codex (GPT-6; exact serving revision not exposed) orchestrated the batch and contributed formatting/repository review. No human mathematical authorship, independent human proof checking, or formal verification is claimed. The historical preparation paragraphs in the reproduced catalogue body concern that body only.

Contribution prompt A was applied to the complete definition and supplied source record, followed by the adversarial checks of prompt B. Local reproduction files are \texttt{batch013/checks.py}, \texttt{computations.json}, and the per-ID audit JSON. They were not public links at preparation time. Numerical checks are finite and do not prove unrestricted results. No LaTeX engine was found locally; compilation was not performed by this research agent.

\end{document}
